% sections/discussion.tex -- what the results mean, limitations, when to use what.
% Every number is a \Num macro (generated/numbers.tex).
\section{Discussion}\label{sec:discussion}

\paragraph{Two regimes.}
When the edits create and destroy only small cycles, as in W2, W4 and the larger
W3 instances, the latency of the DRB methods depends little on the size of the
complex; on the largest W2 and W4 instances \method{mmhm} is
\NumWTwoNThirtyThreeSEightMmhmSpeedupFullId{} and
\NumWFourPlaneNSixteenThousandDTwoMmhmSpeedupFullId{} times faster than static
recomputation. When an edit creates or destroys a cycle that spans the complex,
as on W5 and on W1 whenever the last hole closes or the first one opens, the cost
is linear in the size of that cycle, as \cref{prop:nonlocal} predicts.
Recomputation is also linear, so the DRB keeps at most a constant factor: on W1
at most \NumWOneSphereLFiveMmhmSpeedupFullId, with several variants slower than
recomputation, and on W1 and W5a its tail latency is worse. Which regime a
workload is in is a property of its cycles, not of the size of its edits.

\subsection{Limitations}\label{sec:limitations}

\paragraph{Coefficients and dimension.}
Every result of this paper is over $\Z$. Version~1 discussed integer
coefficients, torsion and the Smith normal form, and described its update as
working over a chosen coefficient ring; nothing in it supports these statements,
and we withdraw them. Ranks over a field do not see torsion, which requires the
Smith normal form over the integers~\cite{munkres1984elements}, and Betti
numbers over $\Z$ can differ from those over other fields: the real projective
plane has Betti numbers $(1,1,1)$ over $\Z$ and $(1,0,0)$ over the
rationals~\cite{hatcher2002algebraic}. We have neither implemented nor tested
other coefficients. The proofs of \cref{sec:theory} hold in every dimension, but
the engine and all experiments are restricted to dimension at most three.

\paragraph{No locality theorem.}
By \cref{prop:nonlocal} the cost of an update is bounded only in
output-sensitive terms (\cref{prop:cost}). Where no spanning cycle is created
these quantities stayed small: on the W2 and W4 instances of E8 the median
number of column additions per update was at most
\NumEEightWFourPlaneNSixteenThousandDTwoDrbIdAddsMed. Nothing guarantees this
for other workloads.

\paragraph{Rebuild policy and the Morse order.}
The growth threshold $\alpha$ and the period $m$ are heuristics, and their effect
depends on the workload (E7); a method that may rebuild in the Morse order also
pays for the index such a rebuild needs (E4). \Cref{lem:apparent} and its
corollaries describe the state right after a rebuild only
(\cref{rem:morse-scope}). Between rebuilds the Morse order is a heuristic whose
measured benefit is real but modest (E8), and for one-shot static recomputation
the collapse never paid for itself (E1).

\paragraph{Measurements.}
All timings come from one virtual machine, single-threaded runs and one compiler
setting. The run-to-run coefficient of variation of the median latency was
\NumNoiseCVMmhmMed\,\% for \method{mmhm} and \NumNoiseCVFullIdMed\,\% for
\method{full-id}, and one run was an outlier attributable to the machine (E1).
The counts of column additions, the footprints and $\nnz(R)+\nnz(V)$ do not
depend on the machine. \method{gudhi} runs through its Python interface and
computes persistence pairs; its times are a reference point for a widely used
library, not the best possible static recomputation. The dynamic reduced basis
stores $V$ and the footprint index, which recomputation does not need; we did not
measure memory. W4 is truncated at dimension $D$, so it measures the algebra on
non-manifold complexes, not the topology of point clouds (\cref{tab:w4-betti}).

\subsection{When to use what}\label{sec:when}

\paragraph{Surfaces with known connectivity:} combinatorial invariants
(\cref{rem:surfaces}). On the largest W1 sphere, \method{euler} was
\NumWOneSphereLSevenEulerSpeedupFullId{} times faster than static recomputation
and faster still than every DRB method.

\paragraph{Barcodes:} zigzag persistence~\cite{carlsson2010zigzag}.
$\beta_k(K_t)$ is the number of $k$-dimensional intervals of the zigzag barcode
that contain~$t$, so zigzag
algorithms~\cite{carlsson2009zigzag,milosavljevic2011zigzag,maria2024discrete,dey2022fast}
answer our query as well, the incremental algorithm of Carlsson, de Silva and
Morozov~\cite{carlsson2009zigzag} after every update. We did not compare
against zigzag implementations.

\paragraph{Connectivity only:} fully dynamic connectivity, which maintains
$\beta_0$ (and, for graphs, $\beta_1=n_1-n_0+\beta_0$) in polylogarithmic
amortised time per edge edit~\cite{holm2001polylog}; on W5a no method that
stores cycle representatives can match this. For complexes in the $3$-sphere
built by insertions only, the algorithm of Delfinado and
Edelsbrunner~\cite{delfinado1995incremental} needs near-linear total time.

\paragraph{Small complexes, or tail latency under global cycles:} static
recomputation with clearing in insertion order (\method{full-id}), not in the
Morse order. It stores no basis, it was faster on average than two DRB variants
on the smallest W3 instance, and it had the lower p99 on every W1 instance and at
every size of W5a. Within the edit sizes we measured ($s\le64$ on W2) it did not
overtake the DRB methods on average; we did not test larger batches.

\paragraph{Edit-local workloads on large general complexes:} \method{mmhm} with
the growth-triggered rebuild. It was \NumWTwoNThirtyThreeSEightMmhmSpeedupFullId{}
times faster than \method{full-id} on the largest W2 instance,
\NumWFourPlaneNSixteenThousandDTwoMmhmSpeedupFullId{} times on the largest W4
instance, and the fastest method at every size of W3, where the basis grows. Its
cost follows the cycles that the edits create or destroy, not the size of the
complex, and it handles non-manifold complexes in which every Betti number
changes.
